\documentclass[10pt,a4paper]{amsart}
\usepackage[margin=19mm]{geometry}
\usepackage{amsmath,amssymb,mathtools}
\usepackage[scr=rsfs,cal=euler]{mathalfa}
\usepackage{xfrac}
\usepackage[unicode,hidelinks,bookmarks=false]{hyperref}
\newcommand{\ie}{i.e.\ }
\newcommand{\cf}{cf.\ }
\newcommand{\resp}{respectively\ }
\newcommand{\Subsection}{\subsection}
\providecommand{\nobreakdash}{\nobreak\hbox{-}}
\setlength{\emergencystretch}{2em}
\multlinegap0pt
\usepackage{amssymb}
\usepackage{xcolor}
\usepackage{tikz}
\newtheorem{lemma}{Lemma}
\newtheorem{corollary}{Corollary}
\newcommand{\N}{N}
\newcommand{\cod}{\operatorname{c}}
\title{Tuza's conjecture for graphs of maximum degree at most seven}
\author{Anish Gupta}
\date{Source: arXiv:2608.06538v1 (CC BY 4.0)}
\begin{document}
\maketitle

\noindent\emph{Source and licence.} Tuza's conjecture for graphs of maximum degree at most seven. Version arXiv:2608.06538v1 (CC BY 4.0), distributed by arXiv under the Creative Commons Attribution 4.0 International licence, http://creativecommons.org/licenses/by/4.0/, as recorded in the arXiv licence metadata for this exact version and shown on the abstract page https://arxiv.org/abs/2608.06538v1. Full licence notice: \texttt{licenses/arxiv-2608.06538-CC-BY-4.0.txt}.\par\smallskip

\noindent\textit{Editorial scope.} Counted unit: the article's lemma lem:template (source0.tex lines 819--826). The article's complete original proof of it is delivered with it (source0.tex lines 828--866, one \texttt{proof} environment of 39 lines). No mathematics is written, repaired or re-derived: the statement and the proof are the article's own lines, in the article's own notation, and no retained line is substituted or re-flowed.  Retained uncounted as the source-local context the proof needs: lines 687--694; lines 712--723; lines 779--786.  Nothing was omitted from any retained span, so the package carries no omission record.  No retained line contains a citation, so the wrapper carries no bibliography and no bibliography entry.  No retained line contains a figure environment, an image-inclusion command or drawing code, so no image is delivered and the package declares no asset dependency.  Editorial formatting only, in addition: an AMS \texttt{amsart} preamble that restates the article's own declarations and macros used by the retained lines, namely the environments \texttt{lemma}, \texttt{corollary} on separate counters, together with the standard packages those lines need.  The retained environments print in this document as Lemma 1, Corollary 1 in order of appearance, whereas the article prints the same items with the numbering of its own full document.  The complete native source of the article as distributed in the arXiv e-print for this exact version is bundled unchanged as \texttt{sources/207007/source0.tex}.
\begin{lemma}[completeness of the local model]\label{lem:localcomplete}
\begin{enumerate}
  \item[(a)] Every triangle of $G$ containing $u$ or $v$ has all three
    vertices in $W$ and all three edges in $E(L)$; hence it is a triangle
    of $L$.
  \item[(b)] Conversely, every triangle of $L$ is a triangle of $G$.
\end{enumerate}
\end{lemma}

\begin{corollary}[local reducibility criterion]\label{cor:localcriterion}
Let $S$ be a set of pairwise edge-disjoint triangles of $L$ and let
$X\subseteq E(L)$ satisfy
\begin{enumerate}
  \item[(i)] $|X|\le 2|S|$;
  \item[(ii)] every triangle of $L$ containing $u$ or $v$ contains an edge of
    $X$;
  \item[(iii)] every $S$-edge with both endpoints outside $\{u,v\}$ lies in
    $X$.
\end{enumerate}
Then $\{u,v\}$ is reducible in $G$, using $S$ and $X$.
\end{corollary}

\begin{lemma}[what $p(R)$ counts]\label{lem:packing}
Let $\mathcal{T}(R)$ be the set of triangles of the forms
\[
  uvx\ (x\in C),\qquad uxy\ (xy\in R),\qquad vxy\ (xy\in R).
\]
Then the maximum size of a pairwise edge-disjoint subfamily of
$\mathcal{T}(R)$ is exactly $p(R)$.
\end{lemma}

\begin{lemma}[template]\label{lem:template}
Let $G$ be $7$-regular and $uv\in E(G)$ with $c=\cod(uv)\in\{5,6\}$; set
$t=3$ if $c=5$ and $t=1$ if $c=6$. If some $R\subseteq E(H)$ satisfies
\begin{equation}\label{eq:template}
  t+|R|+2q(R)\;\le\;2p(R),
\end{equation}
then $\{u,v\}$ is reducible in $G$.
\end{lemma}

\begin{proof}
Choose a minimum vertex cover $Q$ of $H-R$, so $|Q|=q(R)$, and by
Lemma~\ref{lem:packing} choose an edge-disjoint $S\subseteq\mathcal{T}(R)$
with $|S|=p(R)$. Put
\[
  X=\begin{cases}
    \{uv,\,ua,\,vb\}\cup R\cup\{ux,vx : x\in Q\}, & c=5,\ A=\{a\},\ B=\{b\},\\[2pt]
    \{uv\}\cup R\cup\{ux,vx : x\in Q\}, & c=6,\ A=B=\emptyset.
  \end{cases}
\]
All listed edges lie in $E(L)$, and they are pairwise distinct: $uv$ joins the
two hubs, $ua$ and $vb$ join a hub to an exclusive neighbour, the elements of
$R$ join two vertices of $C$, and the spokes $ux,vx$ join a hub to $C$. Hence
$|X|=t+|R|+2q(R)$, and \eqref{eq:template} gives condition (i) of
Corollary~\ref{cor:localcriterion}.

Condition (ii). By Lemma~\ref{lem:localcomplete} it suffices to consider
triangles of $L$ containing $u$ or $v$; take one containing $u$, the case of
$v$ being symmetric. Its other two vertices lie in
$\N_L(u)=\{v\}\cup C\cup A$. The possibilities are:
\begin{itemize}
  \item $uvx$ with $x\in C$: it contains $uv\in X$.
  \item $uxy$ with $x,y\in C$ and $xy\in E(H)$: if $xy\in R$ then $xy\in X$;
    otherwise $xy$ is an edge of $H-R$, so $Q$ covers it and some endpoint,
    say $x$, lies in $Q$, whence $ux\in X$.
  \item $uax$ with $a\in A$, $x\in C$ (only when $c=5$): it contains
    $ua\in X$.
  \item There is no triangle $uab'$ with $a\in A,b'\in B$, since
    $B\cap \N(u)=\emptyset$; and none with two vertices of $A$, since
    $|A|\le1$ here.
\end{itemize}
Condition (iii). The members of $S$ are of the forms $uvx$, $uxy$, $vxy$ with
$xy\in R$. The edges of $uvx$ are $uv,ux,vx$, all incident with $\{u,v\}$. The
only edge of $uxy$ (respectively $vxy$) not incident with $\{u,v\}$ is
$xy\in R\subseteq X$. So every $S$-edge with both endpoints outside $\{u,v\}$
lies in $X$.

Corollary~\ref{cor:localcriterion} now applies.
\end{proof}

\end{document}
